%==============================================================================
% Book: How Our Common Home Works
% A Journey Through the Universe Without Mysteries and Invisible Monsters
% Author: Viktor Logvinovich
%==============================================================================
\documentclass[12pt,a4paper,oneside]{book}

% Encodings and language
\usepackage[T1]{fontenc}
\usepackage[utf8]{inputenc}
\usepackage[english]{babel}
\usepackage{indentfirst}

% Formatting
\usepackage{geometry}
\geometry{margin=2.5cm, top=3cm, bottom=3cm}
\setlength{\headheight}{28pt} % Increased height for long headers
\usepackage{microtype}
\usepackage{titlesec}
\usepackage{fancyhdr}
\usepackage{amsmath, amssymb}
\usepackage{array}
\usepackage{booktabs}
\usepackage{graphicx}
\usepackage{longtable} % Package for breaking long tables across pages

% Hyperlinks (load last)
\usepackage{hyperref}
\hypersetup{colorlinks=true, linkcolor=black, urlcolor=blue}

% Header configuration
\titleformat{\chapter}[display]
  {\normalfont\huge\bfseries\centering}{\chaptertitlename\ \thechapter}{20pt}{\Huge}
\titlespacing*{\chapter}{0pt}{40pt}{30pt}
\titleformat{\section}
  {\normalfont\Large\bfseries}{\thesection}{1em}{}
\titleformat{\subsection}
  {\normalfont\large\bfseries}{\thesubsection}{1em}{}

% Headers and footers
\pagestyle{fancy}
\fancyhf{}
\fancyhead[LE,RO]{\thepage}
\fancyhead[RE]{\leftmark}
\fancyhead[LO]{\rightmark}
\fancyfoot[C]{\small IT3 Paradigm}

%==============================================================================
\begin{document}

%==============================================================================
\begin{titlepage}
    \centering
    \vspace*{2cm}
    {\Huge\bfseries How Our Common Home Works\par}
    \vspace{0.5cm}
    {\Large A Journey Through the Universe Without Mysteries and Invisible Monsters\par}
    \vspace{2cm}
    {\Large Viktor Logvinovich\par}
    \vfill
    {\large Popular Science Edition Based on the IT3 Series\par}
    \vspace{1cm}
    {2026\par}
\end{titlepage}

%==============================================================================
\frontmatter
\tableofcontents

%==============================================================================
\chapter*{Introduction: A Puzzle Missing Some Pieces}
\addcontentsline{toc}{chapter}{Introduction: A Puzzle Missing Some Pieces}

Imagine you are putting together a giant jigsaw puzzle of the night sky. Almost everything is in place, but three pieces are missing right in the center. You look for them under the sofa, in the box, on the floor—they are nowhere to be found. And then you say: “Alright. It means these pieces just don't exist. But to keep the puzzle from falling apart, we'll add some invisible glue and invisible struts.”

Sounds strange, right? But this is exactly what the main cosmological model looks like today. Astronomers see that galaxies spin faster than they should. Instead of rethinking how space itself is structured, we say: “There is invisible \textbf{Dark Matter}.” We see that the Universe is expanding at an accelerating rate. Instead of looking for the cause in the shape of the cosmos, we say: “There is invisible \textbf{Dark Energy} at work.”

We complicate the world, populating it with invisible entities, just to save an old habit of thinking that the Universe is infinite in all directions.

This book offers a different path. The path of simplicity.
What if we throw away all this “dark” clutter? What if we simply imagined the shape of our cosmic home incorrectly?

Welcome to the IT3 paradigm. There is no magic here. There is only geometry, logic, and the amazing harmony of a finite world. Let's go.

%==============================================================================
\mainmatter

%==============================================================================
\chapter{The Shape of Our Home: Finiteness Without Edges}

\section{Infinity or a Room Without Walls?}

We are used to thinking that the Universe is infinite. It's convenient. If the cosmos has no boundaries, light can travel forever without hitting anything. Gravity doesn't have to “loop back.”

But infinity comes at a price. If the Universe is infinite, waves of absolutely any size should exist within it. Imagine an ocean. In an infinite ocean, tsunamis the size of a continent could roam. However, when astronomers look at the cosmic microwave background (the most ancient “echo” of the Big Bang), they see a strange thing: the largest temperature waves are simply not there. Their power is 20\% lower than theory predicts.

Why? Because the waves don't fit into the “room.”

If you stretch a guitar string and pluck it, it cannot produce a wave longer than itself. The lowest bass notes simply won't fit into a short string. The Universe works in exactly the same way. Fluctuations (oscillations) that are larger than its own size cannot exist within it. And data from the Planck satellite shows exactly this: a cutoff at the largest scales. The Universe has a specific size.

When scientists plug this data into equations, the math strictly limits this size. The result is a number: $L_x = 28.57$ billion light-years. This isn't a random figure. It is exactly the scale needed to “cut off” the extra low frequencies in the cosmic microwave background, without breaking everything else.

\section{The “Pac-Man” Effect}

But if the Universe is finite, where is its edge? What happens if you fly all the way to the wall?

There is no wall. Think back to the old arcade game. If your hero runs off the right edge of the screen, he instantly appears on the left side. If he flies up, he pops out from the bottom. The map is closed, but it has no boundaries. Mathematicians call this shape a \textbf{torus} (or a donut).

\textbf{What is a torus in simple terms?} It is a closed space with no edges. You can fly forward forever and eventually return to your starting point. Our Universe is a three-dimensional torus.

This changes everything. In an infinite Universe, energy and gravity dissipate into nowhere. In a finite torus, they cannot disappear. They reflect, overlap upon themselves, and make space “hum” like the walls of a musical instrument.

\section{Why \texorpdfstring{$\sqrt{2}$ and $\sqrt{3}$}{square root of 2 and square root of 3}?}

Here is where the second detail comes into play. If the sides of our “cosmic box” relate as simple numbers (e.g., $1:2:3$), light would reflect along the exact same trajectories, creating artificial patterns and echoes. For the Universe to look random and homogeneous, the sides must relate to each other as \textbf{irrational numbers}: $1 : \sqrt{2} : \sqrt{3}$.

\textbf{What are irrational ratios?} These are proportions that cannot be expressed as a simple fraction. They guarantee that the paths of light and gravitational waves will never loop back perfectly. They fill space evenly, simulating the observed randomness while maintaining finiteness. No mysticism. Just the pure math of ergodicity: deterministic chaos that looks like a perfect mess.

%==============================================================================
\chapter{Cosmic Echo and the Missing Bass}

\section{Cosmic Microwave Background: A Photograph of the Infant Universe}

When the Universe was only 380,000 years old, it cooled down enough for light to travel freely through space. This ancient light is still traveling around us today. We call it the \textbf{Cosmic Microwave Background} (or CMB). It is like a baby photograph of the cosmos.

Tiny spots are visible in this photograph: some places are hotter, some are colder. Scientists break them down into “notes” (multipoles). Low notes (low $\ell$) are the longest, smoothest waves. High notes are short, frequent ripples.

\section{The Low Multipole Anomaly}

The standard model predicts how much energy should be in each “note.” But the lowest, deepest “bass notes” turned out to be quieter than expected. They are missing.

In an infinite Universe, this is a mystery. In our finite one, it is a natural fact. The \textbf{infrared cutoff} (a scientific term for the filtering of long waves) acts as a filter: if a wave does not fit into the size of the torus, it simply does not form. A musical box cannot produce a sound lower than its own length. The Universe demonstrates this clearly.

%==============================================================================
\chapter{The Speedometer Paradox: Why Both Sides Are Right?}

\section{Hubble Tension}

One of the hottest mysteries in modern physics is the discrepancy in measurements of the Universe's expansion rate. This rate is called the \textbf{Hubble constant} ($H_0$).

\begin{itemize}
    \item If we look at the early cosmos (CMB), we get about \textbf{67 km/s per megaparsec}.
    \item If we look at nearby supernovae and Cepheids (beacon stars), we get about \textbf{73}.
\end{itemize}

Is someone making a mistake? Or is the data lying? This discrepancy is called the \textbf{Hubble tension}.

\section{The Solution Through the “Phases” of Our Home}

In the IT3 paradigm, \textbf{both sides are right}. They are simply measuring the Universe at different phases of its maturation.

Early probes see the cosmos before large-scale structures formed. Back then, space was smooth. Late probes see the cosmos after giant voids and filaments of galaxies appeared. At that moment, a special mechanism turns on, accelerating the expansion.

This is not a conflict of data. It is the signature of structure evolution. One Universe, two operating modes. It’s as if you measured the speed of a river in a dry riverbed and during a flood: the readings are different, but it’s the same river.

%==============================================================================
\chapter{The Cosmic Sponge: When Emptiness Starts Working}

\section{Voids: More Than Just Emptiness}

When astronomers map the Universe, they don't see an even fog of stars, but a giant web. Galaxies are gathered into glowing threads, and between them yawn incredible empty spaces — \textbf{voids}. They occupy about \textbf{75\% of the volume of the cosmos}.

It used to be thought that these voids were just “noise,” a byproduct of gravity. But imagine a regular kitchen sponge. It consists of a solid framework and little holes. If you pour water on it, it won't just flow randomly, but through the channels of least resistance.

The cosmos works the same way. Voids are not the absence of something. They are \textbf{channels} through which energy flows completely differently. Together, they form the \textbf{Cosmic Sponge}.

\section{How the Universe Stepped on the Gas (Percolation)}

At the very beginning of time, there were no voids. Matter was distributed almost perfectly evenly. The sponge was “turned off.” But as it expanded, gravity gathered galaxies into clusters, and the voids grew. About 7–8 billion years ago, the voids connected into a single global network. A phase transition occurred, which physicists call \textbf{percolation}.

\textbf{What is percolation?} It is the moment when isolated voids suddenly connect into continuous pathways. Like water finally finding a way through coffee grounds in a coffee maker. Or like foam in a bathtub that suddenly becomes a single mass.

At this moment, the “cosmic sponge” turned on. It began to amplify topological effects almost fourfold. And it is this activation that changes the history of the Universe's expansion in the late epoch.

\section{Topological Sink: Replacing Dark Energy}

When structures form, they release gravitational energy. In an infinite space, it flies off into eternity. In a finite torus, it cannot leave. It reflects, interferes, and focuses in symmetrical points of space. This creates a macroscopic \textbf{negative pressure}. We call this the \textbf{topological sink}.

In simple terms: matter is not just diluted during expansion; part of its energy continuously “flows” into the geometry of space, creating a repulsive effect. When the sponge turned on, this sink intensified. Space began to expand faster not because new “dark” energy was poured into it, but because the conductivity of the fabric of reality itself changed. Acceleration without magic. Just finiteness and a spongy structure.

%==============================================================================
\chapter{Cosmic Friction: Why Don't Galaxies Clump Together?}

\section{S8 Tension}

There is another mystery: the \textbf{S8 tension}. Computer simulations say that galaxies should gather into clusters more actively than we actually see. Reality shows less “clumping.”

\section{Geometric Friction}

In our model, that exact same topological sink adds additional “friction” to the growth of structures. Imagine trying to roll a snowball, but the snow is a bit wet and sticky: the ball grows slower. The same is true here: some of the energy that should have gone into the growth of galaxies escapes into the geometry of the torus through the sponge.

This suppresses structure formation by about 10\%. The result? The theory perfectly matches observations of \textbf{weak gravitational lensing} (when the light of distant galaxies is slightly distorted, like in a funhouse mirror, allowing us to “weigh” invisible mass). Dark matter as a particle is not needed. There is only the curvature of paths in a finite resonator.

%==============================================================================
\chapter{The Music of Particles: Where Does Mass Come From?}

\section{The Higgs Boson and its 125 GeV}

There is a famous mystery in elementary particle physics: why does the Higgs boson weigh exactly 125 GeV? In the standard model, this number is simply written in by hand. It seems random.

But remember that our finite Universe is like a musical instrument. If you take a pipe of a different length, it will sound at different notes. The vacuum in space is not empty. It is filled with quantum fluctuations, which also obey the rules of the resonator. Their energy, and therefore the masses of the particles that interact with them, depend on the size of the “box.”

\section{Emergence: A Property Born Together}

It turns out that the Higgs mass is not “hardcoded” into the particle from the start. It is \textbf{emergent}.

\textbf{What is emergence?} It is a property that appears only when many parts work together, like a wave in a stadium: one person is not a wave, but a thousand people are. The same goes for mass: it arises as a collective effect of the interaction between the topology of space, the percolation network of voids, and the holographic bound of the vacuum.

Mathematics shows a simple connection: if you substitute the size of the torus ($28.57$ Gpc), the amplification factor of the sponge ($\approx 3.8$), and the geometric coefficients of the irrational torus, you get:
\[
m_H \approx 125.1 \text{ GeV}.
\]
The coincidence with the experiment is not random. It means that the Higgs mass is a “note” that our Universe plays thanks to its size and shape.

%==============================================================================
\chapter{How We Test It: Science Doesn't Take Your Word for It}

A good theory is not one that cannot be refuted, but one that \textit{can} be tested. This is called \textbf{falsifiability}. The IT3 paradigm does not hide behind pretty words. It makes clear predictions:

\begin{enumerate}
    \item \textbf{Matching circles in the CMB.} If the Universe is closed, light from the same regions must reach us via different routes. Pairs of circles with identical temperature structures should remain on the cosmic microwave background map. IT3 predicts their angular radius: $30^\circ$–$60^\circ$. Future telescopes (CMB-S4, LiteBIRD) will be looking for them.
    
    \item \textbf{Evolution of the Hubble constant.} If the sponge turns on gradually, then the derived $H_0$ should smoothly increase at redshifts $z < 1$. The DESI and Euclid surveys will measure this to within a percent.
    
    \item \textbf{Suppression of gravitational waves.} Inflation predicts a noticeable background of primordial waves. In a finite torus, long waves simply do not fit. IT3 predicts that they should be almost non-existent ($r < 0.01$).
    
    \item \textbf{Anomalies in supervoids.} The flow of energy through percolating voids should leave a mark in the CMB. Its amplitude should be about twice as large as in the old model.
\end{enumerate}

If these signatures are not found, the paradigm will be refined or replaced. That is how science works. But while the stars are silent, the equations have already come together into a self-consistent picture.

%==============================================================================
\backmatter

%==============================================================================
\chapter*{Conclusion: One Home, One Geometry}
\addcontentsline{toc}{chapter}{Conclusion: One Home, One Geometry}

For decades, modern cosmology has added invisible components to maintain an assumption: space is infinite. The IT3 paradigm asks a simpler question: \emph{what if space is finite, and we simply forgot to account for how energy flows in a closed shape?}

The answer is surprisingly consistent:
\begin{itemize}
    \item \textbf{Shape} cuts off the longest waves, explaining the low multipole anomaly.
    \item \textbf{Flow} diverts energy into topological antinodes, replacing dark energy with a physical dissipation mechanism.
    \item \textbf{Structure} channels this flow through voids, activating acceleration and suppressing excess clustering.
    \item \textbf{Resonance} of the vacuum in a finite torus sets the mass of fundamental particles.
\end{itemize}

No new particles. No fine-tuning. No multiverse. Just the geometry of a finite torus, the thermodynamics of dissipation, and the observed spongy distribution of matter. One parameter ($L_x$) instead of a dozen arbitrary tweaks.

The next generation of surveys will decisively test these predictions. If the data aligns, cosmology will transition from a paradigm of invisible substances to a paradigm of visible geometry. If not, the model will be refined. This is not a defeat. It is progress.

But for now, the theoretical framework is complete, self-consistent, and ready for nature to deliver its verdict. The Universe might not be as mysterious as we thought. Perhaps it just needed the right shape for everything to fall into place.

\vspace{2cm}
\begin{flushright}
\textit{“While scientists search for circles, geometry is already working.”}
\end{flushright}

%==============================================================================
\chapter*{Glossary for Young Explorers}
\addcontentsline{toc}{chapter}{Glossary for Young Explorers}

A quick reference guide for all scientific terms used in the book.

\renewcommand{\arraystretch}{1.3}
\begin{longtable}{>{\bfseries}p{4.5cm}p{10cm}}
\toprule
\textbf{Term} & \textbf{Simple Explanation} \\
\midrule
\endfirsthead

\multicolumn{2}{c}{{\small\textit{Glossary (continued)}}}\\
\toprule
\textbf{Term} & \textbf{Simple Explanation} \\
\midrule
\endhead

\bottomrule
\endfoot

\bottomrule
\endlastfoot

\textbf{Cosmology} & The science of the Universe as a whole: how it originated, how it is structured, and how it will develop. \\
\textbf{Topology} & A branch of mathematics that studies how space is connected to itself (e.g., whether it is closed, whether it has “holes” or edges). \\
\textbf{Torus (3-torus)} & A closed space without edges. Like the screen of an old video game: you go off the right edge and appear on the left. \\
\textbf{Irrational ratios} & Proportions of the sides of a space that cannot be expressed as a simple fraction (e.g., $1:\sqrt{2}:\sqrt{3}$). They guarantee that light and waves will never loop identically, creating apparent randomness. \\
\textbf{Cosmic Microwave Background (CMB)} & The most ancient light in the Universe, the “echo” of the Big Bang. It reaches us from everywhere and shows what the Universe was like in its infancy. \\
\textbf{Low multipoles (low $\ell$)} & The longest, smoothest temperature waves in the CMB. The equivalents of “bass notes” in a cosmic symphony. \\
\textbf{Infrared cutoff} & A natural limitation: waves longer than the space itself cannot exist in a finite space. The long “basses” simply do not fit. \\
\textbf{Hubble Constant ($H_0$)} & The speed at which the Universe is expanding right now. It is measured in kilometers per second for every megaparsec of distance. \\
\textbf{Hubble Tension} & A disagreement between two ways of measuring the expansion rate: via the early cosmos ($\sim 67$) and via nearby objects ($\sim 73$). In IT3, this is not an error, but a sign of the Universe changing its “phase.” \\
\textbf{Voids} & Giant cosmic empty spaces where there are almost no galaxies. They occupy $\sim 75\%$ of the volume of the Universe. \\
\textbf{Cosmic sponge} & The large-scale structure of the Universe, where galaxies form threads and voids form empty spaces. It acts like a waveguide, amplifying energy flows. \\
\textbf{Percolation} & The moment when isolated empty spaces connect into a single network. Like water finding a way through coffee grounds. In space, this “turns on” accelerated expansion. \\
\textbf{Topological sink} & A mechanism in which energy in a closed space does not dissipate into infinity, but is focused by geometry itself, creating negative pressure and repulsion. It replaces “dark energy.” \\
\textbf{Dark Energy / Dark Matter} & Hypothetical invisible components that were added “by hand” in the old model to explain the acceleration and rotation of galaxies. In IT3, their role is played by geometry and the structure of space. \\
\textbf{S8 Tension} & An observed discrepancy: there are fewer galaxies in clusters than simulations predict. In IT3, this is explained by “cosmic friction” from the topological sink. \\
\textbf{Weak gravitational lensing} & The bending of light from distant galaxies caused by gravity. It allows us to “weigh” invisible mass by how it bends light. \\
\textbf{Higgs Boson} & An elementary particle associated with a field that gives mass to other particles. \\
\textbf{Emergence} & A property of a system that only arises when many parts interact (just as temperature arises from the movement of molecules). In IT3, mass is not “hardcoded” into particles, but is born from the resonance of space. \\
\textbf{Falsifiability} & The property of a scientific theory to be testable and potentially disproven by observations. If the predictions do not come true, the theory will be changed or discarded. \\
\textbf{Matching circles} & Expected patterns in the CMB: if the Universe is closed, light travels around it via different paths, and paired circles with identical temperatures should appear on the sky. \\
\end{longtable}

\vspace{2cm}
\begin{flushright}
\small April 2026 | Kobrin, Belarus\\
\small Code and data repository: \url{https://github.com/Viktar-Pi/FlatIrrationalTorus}
\end{flushright}

%==============================================================================
\end{document}
